% !PRINTED 381
\weberpage{381}
%% [running head, not body text -- commented out 2026-08-07] Page furniture (running
%% head, folio, signature) is excluded from this edition by standing policy.
%% §\,104. Characters and genera of the quadratic forms. \hfill 381

If $\delta_1, \delta_2$ have no common divisor, then this symbolic
multiplication is an actual one. But if they do have a common
divisor, then a square factor has still to be cast out of the true
product $\delta_1\delta_2$. If we set $\delta_1 = \delta_2$ in (2),
then $\delta = 1$, that is, every element of the group of the
$\delta$ is reciprocal to itself, and together with (2) there holds:

\weq{3}{\delta_1 = \delta\delta_2.}

The elements of this group can be represented in the following way by a
\gesperrt{basis} (Bd. II, § 11):

\weq{4}{\delta = \delta_1^{\varepsilon_1}\delta_2^{\varepsilon_2}\dots\delta_\lambda^{\varepsilon_\lambda},}

in which $\varepsilon_1, \varepsilon_2, \dots \varepsilon_\lambda$ have the values 0 or 1, and $\lambda$ has the
meaning given above. $\delta_1, \delta_2, \dots \delta_\lambda$ are, in case 1), the odd
prime discriminants $\pm p$ that divide $D$; in case 2) there
is added to these $-4$, in case 3) $8$, in case 4) $-8$, in
case 5) $8$ and $-8$ (or $-4$ and $8$ or $-4$ and $-8$).

With the help of the fundamental divisors the \gesperrt{characters} of the
primitive quadratic forms are now defined.

In a class of forms $A$ we choose a representative
$(a,b,c)$ in which $a$ is positive and relatively prime to $D$. If then a
number $m$ coprime to $D$ is representable by $A$, we have
\[
m = ax^2 + bxy + cy^2,
\]
\weq{5}{4am = (2ax + by)^2 - Dy^2.}

Now let $\delta$ be a fundamental divisor of $D$ and $m$ \gesperrt{relatively prime}
to $\delta$; then, first for an odd $\delta$, by § 85, (14), (16):

\weq{6}{(\delta, m) = (\delta, a).}

But this relation holds also for $\delta = -4\delta_1, +8\delta_1, -8\delta_1$.
For in these cases $a$ and $m$ are odd, and by (1)
and (5):
\[
\begin{aligned}
\delta &= -4\delta_1, & D &\equiv 0, -4 \ (\mathrm{mod}\ 16), & am &\equiv 1 \ (\mathrm{mod}\ 4), \\
\delta &= 8\delta_1, & D &\equiv 0, 8 \ (\mathrm{mod}\ 32), & am &\equiv \pm 1 \ (\mathrm{mod}\ 8), \\
\delta &= -8\delta_1, & D &\equiv 0, -8 \ (\mathrm{mod}\ 32), & am &\equiv 1, 3 \ (\mathrm{mod}\ 8),
\end{aligned}
\]
from which, by § 85, equation (6) follows for these cases also,
and it is thereby proved in general.

4. \gesperrt{The value of the symbol $(\delta, m)$ therefore depends not on
the number $m$, but only on the class of forms $A$
by which $m$ is representable. It is called the character of this class
belonging to the fundamental divisor $\delta$, and is denoted
by $\chi(\delta, A)$.}
